%\documentclass[spanish]{beamer}

\documentclass{beamer}
\usepackage{pgf,pgfpages}
%\pgfpagesuselayout{4 on 1}[letterpaper,landscape,border shrink=0.5in]
\usepackage{algorithm,algorithmic}
%\usepackage{algorithm,algorithmicx}
%\usepackage[noend]{algpseudocode}
\usepackage{listings}
\lstset{basicstyle=\ttfamily,
	showstringspaces=false,
	commentstyle=\color{red},
	literate={á}{{\'a}}1
	{ã}{{\~a}}1
	{é}{{\'e}}1
	{É}{{\'E}}1
	{ó}{{\'o}}1
	{í}{{\'i}}1
	{ñ}{{\~n}}1
	{¡}{{!`}}1
	{¿}{{?`}}1
	{ú}{{\'u}}1
	{Í}{{\'I}}1
	{Ó}{{\'O}}1
}

\usepackage{beamerthemesplit}
\usepackage{color}
\usepackage[utf8]{inputenc}
\usepackage[spanish]{babel}
%\usepackage[latin1]{inputenc}
%\usepackage[pdftex]{graphicx}
%\usepackage{algorithm}
%\usepackage{algorithmic-C}
%\usepackage{algorithmic}
\usepackage{amsthm}
\usepackage{multicol}
\usepackage{multirow}
\usepackage{fancyvrb,relsize}
\usepackage{hhline}
\usepackage{tikz}
\usetikzlibrary{trees, shapes, calc}
\newcommand{\inputTikZ}[2]{%
	\scalebox{#1}{\input{#2}}
}

\usetheme{Warsaw} %Gueno
%\usetheme{Darmstadt} %OK!
% \usetheme{Frankfurt} %OK!
%\usetheme{Goettingen}
%\usetheme{Dresden}
%\usetheme{JuanLesPins} %OK!!
%\usetheme{Marburg}
%\usetheme{Montpellier}
%\usetheme{Rochester} %sobrio
%\usetheme{Singapore}
%\usetheme{Szeged}

%\usecolortheme{albatross}
%\usecolortheme{seahorse}
%\usecolortheme{beetle}
%\usecolortheme{crane}
%\usecolortheme{dolphin}
%\usecolortheme{dove}
%\usecolortheme{fly}
%\usecolortheme{lily} %Gueno
\usecolortheme{orchid}%Gueno
%\usecolortheme{rose}
%\usecolortheme{seagull}
%\usecolortheme{whale}


\DefineVerbatimEnvironment%
      {VerbatimNN}%
      {Verbatim}%
      {fontsize=\relsize{-1}}
\DefineVerbatimEnvironment%
      {VerbatimNNN}%
      {Verbatim}%
      {fontsize=\relsize{-2}}

\DefineVerbatimEnvironment%
      {VerbatimPseudo}%
      {Verbatim}%
      {numbers=left, fontsize=\relsize{-1}}
\DefineVerbatimEnvironment%
      {VerbatimPseudoReduced}%
      {Verbatim}%
      {numbers=left, fontsize=\relsize{-2}}
\DefineVerbatimEnvironment%
      {VerbatimPseudoReduced2}%
      {Verbatim}%
      {numbers=left, fontsize=\relsize{-3}}
\DefineVerbatimEnvironment%
      {VerbatimC}%
      {Verbatim}%
      {commandchars=@ªº, numbers=left, fontsize=\relsize{-1}}
\DefineVerbatimEnvironment%
      {VerbatimReducedC}%
      {Verbatim}%
      {commandchars=@ªº, numbers=left, fontsize=\relsize{-2}}
\DefineVerbatimEnvironment%
      {VerbatimReduced2C}%
      {Verbatim}%
      {commandchars=@ªº, numbers=left, fontsize=\relsize{-3}}



%\AtBeginSubsection[]
%{
% \begin{frame}
% \frametitle{Sinopsis}
% Aqui decides: cada seccion y subseccion:
% \tableofcontents[currentsection,currentsubsection]
 %O solamente la seccion (Solo uno de ellos)
%\tableofcontents[currentsection]
% \end{frame}
%}

% Si quieres que todo por default aparezca de poco en poco
% descomenta esta linea
%\beamerdefaultoverlayspecification{<+->}

\setbeamertemplate{footline}{\hfill\insertframenumber/\inserttotalframenumber}
\setbeamertemplate{navigation symbols}{}

\setbeamercovered{transparent}

\title{ELO320 - Estructuras de Datos y Algoritmos\\Introducción a C\\Struct y Strings}
%\subtitle{Ingeniería Civil Telemática UTFSM}
\author{Dr. Nicolás Gálvez R.\\{\small nicolas.galvez@usm.cl}}
\institute{Ing. Civil Telemática - Departamento de Electrónica}
\titlegraphic{
\includegraphics[scale=.18]{images/logoutfsm}}
\date{2do Semestre 2022}

\begin{document}

\frame{\titlepage}

%\section*{Sinopsis}
%\frame{  \frametitle{Sinopsis}\tableofcontents[pausesections]}
%
%
\section{Struct: Agrupando información.}
\begin{frame}[fragile]{Estructuras}
\begin{itemize}
	\item Conjunto heterogéneo de elementos de datos.
	\item Variable: campo, miembro o atributo.
	\item Base de los Tipos de Datos Abstractos (TDA).
	\item Encapsulamiento de información.
	\item Acceso: Operador punto (\texttt{.}).
\end{itemize}
\begin{exampleblock}{C}
\begin{lstlisting}[language=C, basicstyle=\tiny, escapeinside={|}{|}]
struct empleado {
    char nombre[40];
    float sueldo;
};

int main(){
	struct empleado e1;
	e1.sueldo = 550000.00; //acceso a sueldo
	strcpy(e1.nombre, "Juan"); //acceso a nombre
	printf("El empleado |\%|s gana |\%|.2f USD al año", e1.nombre, e1.sueldo); 
	return 0;
}
\end{lstlisting}
\end{exampleblock}
\end{frame}

\begin{frame}[fragile]{Estructuras y punteros.}
\begin{itemize}
	\item struct $\rightarrow$ nuevo tipo de datos $\rightarrow$ punteros.
	\item El operador punto se combina con derreferencia, e.g., \texttt{(*var).nombre};
	\item Ojo, \texttt{(*var).nombre $\neq$ *var.nombre.}
	\item En vez, se puede usar el operador \verb+->+. 
\end{itemize}
\begin{exampleblock}{C}
\begin{lstlisting}[language=C, basicstyle=\tiny, escapeinside={|}{|}]
struct empleado {
    char nombre[40];
    float sueldo;
};

int main(){
	struct empleado e1;
	struct empleado *ptr;
	ptr = &e1;
	ptr->sueldo = 550000.00; //acceso a sueldo
	strcpy(ptr->nombre, "Juan"); //acceso a nombre
	printf("El empleado |\%|s gana |\%|.2f USD al año", ptr->nombre, ptr->sueldo); 
	return 0;
}
\end{lstlisting}
\end{exampleblock}
\end{frame}

\begin{frame}[fragile]{Definiendo nuevos tipos de datos}
\begin{itemize}
	\item \textbf{typedef} $\rightarrow$ crear nuevos tipos de datos. 
	\item Se utiliza para encapsular structs. 
	\item También para no trabajar con tipos de datos comunes de forma directa.
\end{itemize}
\begin{exampleblock}{C}
\begin{lstlisting}[language=C, basicstyle=\tiny, escapeinside={|}{|}]
typedef float real;

typedef struct empleado {
    char nombre[40];
    real sueldo;
} empleado_t;

int main(){
	empleado_t e1;
	e1.sueldo = 550000.00; //acceso a sueldo
	strcpy(e1.nombre, "Juan"); //acceso a nombre
	printf("El empleado |\%|s gana |\%|.2f USD al año", e1.nombre, e1.sueldo); 
	return 0;
}
\end{lstlisting}
\end{exampleblock}
\end{frame}


\section{Strings}
\begin{frame}[fragile]{Strings}
\begin{itemize}
	\item Estructura de Datos para poder representar palabras.
	\item Trabajan sobre arreglos de caracteres.
	\item La norma es que finalizen con el caracter fin de string: '\textbackslash 0'.
	\item \textbf{string.h} $\rightarrow$ funciones.
	\item Dependen de la existencia del caracter '\textbackslash 0'.
	\item Indexados como todo arreglo.
	\item Comillas (") $\rightarrow$ string, apóstrofe (') $\rightarrow$ char.
 	\item TDA String: funciones + arreglo de caracteres + '\textbackslash 0'.
\end{itemize}

\begin{center}
\includegraphics[width=.8\linewidth]{images/string}
\end{center}
\end{frame}

\begin{frame}[fragile]{Definición}
\begin{exampleblock}{C}
\begin{lstlisting}[language=C, basicstyle=\tiny, escapeinside={|}{|}]
int main(){

	char a[10]; //Tamaño fijo.
	char *b; // Tamaño dinámico.


	a="Hola"; // inválido
	b="Hola"; // válido

	return 0;
}

\end{lstlisting}
\end{exampleblock}
\end{frame}

\begin{frame}[fragile]{Inicialización}
\begin{exampleblock}{C}
\begin{lstlisting}[language=C, basicstyle=\tiny, escapeinside={|}{|}]
int main(){

	char a[10]="Hola"; //Tamaño fijo con valor inicial.
	char a2[10]={'H','o','l','a','\0'}; //obligación el '\0'

	char b[5]="Juan"; // Éstas dos son equivalentes.
	char b2[]="Juan" // Tamaño fijo definido con valor inicial.
	
	char *c = "ELO320"; // String constante, tamaño variable. 

	char *d = malloc(...); // String dinámico, puntero.
	strcpy(d, "String Pro");
	free(d); //próx. clase. 	

	a[0]='D'; //válido
	b[0]='D'; //válido
	c[0]='D'; //no valido, segmentation fault.
	d[0]='D'; //valido, por algo es pro.

	return 0;
}

\end{lstlisting}
\end{exampleblock}
\end{frame}

\begin{frame}[fragile]{Biblioteca string.h}
\begin{itemize}
	\item \verb+#include<string.h>+
	\item Contiene funciones y elementos útiles para manejar strings.
\end{itemize}

\begin{exampleblock}{C}
\begin{lstlisting}[language=C, basicstyle=\tiny, escapeinside={|}{|}]
int main(){

	char a[10];
	char *b, *c;
	int tam;

	strcpy(a,"Hola"); // strcpy(a,b) copia en a lo de b
	//strncpy(a,b,n) copia los n primeros caracteres.
	//retornan un puntero a "a". 

	b=" Mundo";
	
	c=strcat(a,b); // concatena a con b;
	//strncat concatena a "a" los primeros n caracteres de b.
	//retornan un puntero a "a". 
	
	tam=strlen(a); //largo de a. no cuenta '\0'.
	// retorna size_t, es un typedef de unsigned long int.

	printf("tamaño: |\%|d\n", tam);

	return 0;
}

\end{lstlisting}
\end{exampleblock}
\end{frame}


\begin{frame}[fragile]{Biblioteca string.h}
Mas info:

{\centering\large \url{https://en.wikibooks.org/wiki/C_Programming/string.h}}

\end{frame}

\begin{frame}[fragile]{Desafíos}
\begin{itemize}
\item Escriba un programa que pide al usuario un número entre 1 y 100, y a continuación llama a una función que devuelve la suma de todos los números pares que van del 1 al número que dio el usuario. Al final esta suma se imprime en la pantalla.

\item Escriba un programa que, dada la hora en su computadora, llama a una función que calcula y devuelve el ángulo entre las manecillas del reloj (hora(s) y minutos). Este valor el programa lo imprime en la pantalla. (miren en los códigos de ayuda sobre cómo usar la librería time.h para acceder a la hora en su sistema).

\item Escriba un programa que, dado un arreglo de enteros de dos dimensiones (array [][]) inicializado de manera random (fíjense en las funciones en stdlib.h), definido global, desarrolle un algoritmo que encuentra el mínimo valor en el arreglo entero pero pasando una sola vez por cada elemento del arreglo.
\end{itemize}
\end{frame}


\end{document}
